EEG measurements capture electrical potentials at the scalp surface. Understanding the electromagnetic basis of these potential fields is essential for designing phantoms and simulation models that accurately replicate signal propagation in conductive media.

\subsubsection{Maxwell's Equations}

The behavior of electric and magnetic fields in space and time is governed by Maxwell's equations, which constitute the fundamental laws of classical electromagnetism. In their differential form, these equations are:

\begin{align}
\nabla \cdot \mathbf{D} &= \rho_f, \label{eq:gauss_electric} \\[6pt]
\nabla \cdot \mathbf{B} &= 0, \label{eq:gauss_magnetic} \\[6pt]
\nabla \times \mathbf{E} &= -\frac{\partial \mathbf{B}}{\partial t}, \label{eq:faraday} \\[6pt]
\nabla \times \mathbf{H} &= \mathbf{J} + \frac{\partial \mathbf{D}}{\partial t}.
\label{eq:ampere_maxwell}
\end{align}

Here, $\mathbf{E}$ is the electric field (V/m), $\mathbf{D}$ is the electric displacement field (C/m²), $\mathbf{B}$ is the magnetic flux density (T), and $\mathbf{H}$ is the magnetic field intensity (A/m). The quantity $\rho_f$ denotes the free charge density (C/m³), $\mathbf{J}$ is the current density (A/m²), and $\partial \mathbf{D} / \partial t$ represents the displacement current density (A/m²).

To close this system, constitutive relations connect these field quantities. For linear, isotropic materials, these simplify to:

\[
\mathbf{D} = \varepsilon\,\mathbf{E}, \qquad
\mathbf{B} = \mu\,\mathbf{H}, \qquad
\mathbf{J} = \sigma\,\mathbf{E},
\]

where $\varepsilon$ is the electric permittivity, $\mu$ is the magnetic permeability, and $\sigma$ is the electrical conductivity~\cite{feynman_2011,weik-2000,jackson-1999,griffiths-2017}.

\subsubsection{Quasi-Static Approximation}

At sufficiently low frequencies in spatially compact conductive systems, electromagnetic phenomena can be accurately described using the quasi-static approximation (QSA). This approximation applies when temporal and spatial variations of electromagnetic fields are slow enough that wave propagation effects may be neglected~\cite{plonsey-1967,haus-melcher-1989,jackson-1999}.

The quasi-static regime is characterized by two key conditions that must be simultaneously satisfied for the approximation to be valid~\cite{haus-melcher-1989}.

\paragraph*{Spatial compactness criterion.}
The first condition requires that the characteristic system length scale $L$ be much smaller than the electromagnetic wavelength:

\[
L \ll \lambda = \frac{1}{f\sqrt{\mu\varepsilon}},
\]

where $f$ is the frequency, and $\mu$ and $\varepsilon$ are the magnetic permeability and electric permittivity of the medium, respectively. When this condition holds, spatial phase variations across the domain are negligible, and electromagnetic fields adjust effectively instantaneously throughout the volume~\cite{haus-melcher-1989,jackson-1999}. Under these conditions, retardation and radiation effects can be ignored, and the system lies deep within the quasi-static regime.

\paragraph*{Negligible displacement currents.}
In Ampère's law (Eq.~\ref{eq:ampere_maxwell}), the total current density consists of a conduction component $\mathbf{J}$ and a displacement current component $\partial \mathbf{D}/\partial t$. The relative importance of these two contributions is quantified by the dimensionless ratio:

\[
\frac{|\partial \mathbf{D}/\partial t|}{|\mathbf{J}|} = \frac{\omega\varepsilon}{\sigma},
\]

where $\omega = 2\pi f$ is the angular frequency and $\sigma$ is the electrical conductivity. When

\[
\omega\varepsilon \ll \sigma,
\]

displacement currents are negligible compared to conduction currents, and the electromagnetic response is dominated by Ohmic conduction~\cite{plonsey-1967,haus-melcher-1989}. In this electroquasistatic regime, Ampère's law simplifies to:

\[
\nabla \times \mathbf{H} \approx \mathbf{J}.
\]

Under these conditions, capacitive or inductive energy storage effects become negligible, and the material response is effectively resistive with no appreciable phase lag between the electric field and current density~\cite{haus-melcher-1989}.

\paragraph*{Irrotational electric field.}
Faraday's law (Eq.~\ref{eq:faraday}) also simplifies in the quasi-static limit. Because the currents vary slowly and produce negligible time-varying magnetic flux, the electromagnetic induction term vanishes to a very high degree of accuracy:

\[
\nabla \times \mathbf{E} \approx \mathbf{0}.
\]

As a consequence, the electric field is irrotational and can be expressed as the gradient of a scalar potential:

\[
\mathbf{E} = -\nabla u,
\]

where $u(\mathbf{x},t)$ is the electric potential (V). This relationship is fundamental to modeling potential-based measurements, as it allows the complex electromagnetic problem to be reduced to solving for a scalar potential field that represents the measurable voltages~\cite{plonsey-1967,jackson-1999}.

\subsubsection{Application to bioelectric systems.}
In the context of neuromodulation and bioelectric measurements, the validity of the electro-quasistatic approximation is not merely a heuristic assumption
but has been explicitly formalized and justified in contemporary modeling literature. In particular, \citet{wang-2024}
provide a rigorous definition of the quasistatic approximation as it is applied in neuromodulation pipelines, emphasizing that QSA is invoked
specifically at the stage where electric fields and potentials in tissue are computed from prescribed stimulation sources. Under this framework,
the approximation rests on four assumptions: negligible wave propagation and self-induction in tissue, linear material response, purely resistive behavior,
 and non-dispersive electrical properties over the relevant frequency range.